%% Custom commands
%% ===============

% Alternative to `gamechange` from the cryptocode lib
% need `xcolor` package for that!

\newcommand{\highlight}[1]{%
  \smash{%
    \tikz[baseline=(X.base)]{
      \node[
        fill=gray!20,
        inner xsep=2.5pt,
        inner ysep=1.2pt,
        anchor=base
      ] (X) {$\displaystyle #1$};
    }%
  }%
}

%% Special characters for number sets, e.g. real or complex numbers.
\newcommand{\C}{\ensuremath{\mathbb{C}}}
\newcommand{\K}{\ensuremath{\mathbb{K}}}
\newcommand{\N}{\ensuremath{\mathbb{N}}}
\newcommand{\Q}{\ensuremath{\mathbb{Q}}}
\newcommand{\R}{\ensuremath{\mathbb{R}}}
\newcommand{\Z}{\ensuremath{\mathbb{Z}}}
\newcommand{\F}{\ensuremath{\mathbb{F}}}
\newcommand{\I}{\ensuremath{\mathbb{I}}}

%% Most common use

% Annoying that this is necessary
\AtBeginDocument{ 
  \renewcommand{\O}{\ensuremath{\mathcal{O}}}
}

\newcommand{\A}{\ensuremath{\mathcal{A}}}
\newcommand{\B}{\ensuremath{\mathcal{B}}}
\newcommand{\G}{\ensuremath{\mathcal{G}}}
\renewcommand{\S}{\ensuremath{\mathcal{S}}}
\renewcommand{\K}{\ensuremath{\mathcal{K}}}
\newcommand{\M}{\ensuremath{\mathcal{M}}}
\renewcommand{\L}{\ensuremath{\mathcal{L}}}
\newcommand{\T}{\ensuremath{\mathcal{T}}}
\newcommand{\U}{\ensuremath{\mathcal{U}}}
\renewcommand{\P}{\ensuremath{\mathcal{P}}}
\newcommand{\X}{\ensuremath{\mathcal{X}}}
\newcommand{\Y}{\ensuremath{\mathcal{Y}}}
\newcommand{\V}{\ensuremath{\mathcal{V}}}
\newcommand{\E}{\ensuremath{\mathcal{E}}}

\renewcommand{\a}{\alpha}
\renewcommand{\b}{\beta}
\renewcommand{\c}{\gamma}
\renewcommand{\d}{\delta}
\newcommand{\e}{\varepsilon}
\newcommand{\f}{\varphi}
\newcommand{\g}{\gamma}
\newcommand{\h}{\eta}
\renewcommand{\k}{\kappa}
\renewcommand{\l}{\lambda}
\newcommand{\m}{\mu}
\newcommand{\n}{\nu}
\newcommand{\p}{\pi}
\renewcommand{\r}{\rho}
\newcommand{\s}{\sigma}
\renewcommand{\t}{\tau}
\renewcommand{\u}{\upsilon}
\newcommand{\x}{\ensuremath{\mathbf{x}}}
\renewcommand{\v}{\ensuremath{\mathbf{v}}}
\newcommand{\z}{\ensuremath{\mathbf{z}}}
\newcommand{\y}{\ensuremath{\mathbf{y}}}
\newcommand{\w}{\omega}

% Very natural
\renewcommand{\o}{\circ}

\newcommand{\closure}{\textsf{cl}}
\newcommand{\trace}[1]{\ensuremath{\textsf{trace}(#1)}}

%% Additional simple symbols 
\newcommand{\sample}{\ensuremath{\gets_\$}} 
\newcommand{\bin}{\ensuremath{\{0,1\}}}      
\newcommand{\vrf}{v}                          % for vrf instances
\newcommand{\VRF}{\textup{\textbf{VRF}}}      % for the VRF primitive
\newcommand{\PPTM}{\textup{\textsf{ PPTM }}}
\newcommand{\iddots}{\ensuremath{{\cdot^{\cdot^{\cdot}}}}}
\newcommand{\SK}{\mathcal{SK}}
\newcommand{\PK}{\mathcal{PK}}

\newcommand{\PSPACE}{\ensuremath{\textup{\textsf{PSPACE}}}}
\newcommand{\NP}{\ensuremath{\textup{\textsf{NP}}}}
\newcommand{\PP}{\ensuremath{\textup{\textsf{P}}}}

\newcommand{\err}{\ensuremath{\eps}}
\newcommand{\runtime}{\ensuremath{\mathfrak{t}}}
\newcommand{\errtime}{\ensuremath{(\runtime, \eps)}}
\newcommand{\errtimelen}{\ensuremath{(\runtime, \eps, \mathfrak{a}, \mathfrak{b})}}
\newcommand{\sk}{\ensuremath{\textup{\textsf{sk}}}}
\newcommand{\pk}{\ensuremath{\textup{\textsf{pk}}}}
\newcommand{\pcreq}{\operatorname{\textbf{req}}}
\newcommand{\interp}[1]{\ensuremath{[\![#1]\!]}}
\newcommand{\xor}{\ensuremath{\oplus}}
\newcommand{\rank}[2]{\ensuremath{\text{rk}}_{#1}\Paren*{#2}}

\DeclareMathOperator{\pptm}{p.p.t.m.}
\DeclareMathOperator{\ppt}{ppt}

%% Fixed/scaling delimiter examples (see mathtools documentation)
\let\abs\relax
\let\norm\relax
\DeclarePairedDelimiter\abs{\lvert}{\rvert}
\DeclarePairedDelimiter\norm{\lVert}{\rVert}

% Define generic paired delimiters:
\DeclarePairedDelimiter{\Paren}{(}{)}
\DeclarePairedDelimiter{\Bracket}{[}{]}
\DeclarePairedDelimiter{\Brace}{\{}{\}}

% \begin{figure}[tb]
% 	\begin{pcvstack}[center, space=1em]
% 		\procedure[linenumbering]%
% 		{Experiment $\Exp{VF}{up}(\lambda, \A)$}
% 		{
% 			... 
% 		}
% 	\end{pcvstack}
% 	\caption{Unique provability security experiment.}
% 	\label{fig:up_experiment}
% \end{figure}

% Big-O
\NewDocumentCommand{\Oh}{som}{%
  \ensuremath{\mathcal{O}\IfBooleanTF{#1}{%
    \Paren*{#3}%          % *-version => auto-size
  }{%
    \IfNoValueTF{#2}{%
      \Paren{#3}%         % no optional => normal size
    }{%
      \Paren[#2]{#3}%     % optional => \big, \Big, etc.
    }%
  }%
}}

%Big Omega:
\NewDocumentCommand{\Om}{som}{%
\ensuremath{\Omega\IfBooleanTF{#1}{%
    \Paren*{#3}%
  }{%
    \IfNoValueTF{#2}{%
      \Paren{#3}%
    }{%
      \Paren[#2]{#3}%
    }%
  }%
}}

% Big Theta:
\NewDocumentCommand{\Th}{som}{%
\ensuremath{\Theta\IfBooleanTF{#1}{%
    \Paren*{#3}%
  }{%
    \IfNoValueTF{#2}{%
      \Paren{#3}%
    }{%
      \Paren[#2]{#3}%
    }%
  }%
}}

% little-o:
\NewDocumentCommand{\oh}{som}{%
\ensuremath{o\IfBooleanTF{#1}{%
    \Paren*{#3}%
  }{%
    \IfNoValueTF{#2}{%
      \Paren{#3}%
    }{%
      \Paren[#2]{#3}%
    }%
  }%
}}

% little-omega:
\NewDocumentCommand{\om}{som}{%
\ensuremath{\omega\IfBooleanTF{#1}{%
    \Paren*{#3}%
  }{%
    \IfNoValueTF{#2}{%
      \Paren{#3}%
    }{%
      \Paren[#2]{#3}%
    }%
  }%
}}

% Probability:  \Prob{Event} =>  P{Event},   \Prob*{...} => auto-sized { ... }
%    By default, we use curly braces for probabilities.
\NewDocumentCommand{\Prob}{som}{%
\ensuremath{\Pr% or \mathbb{P}, if you prefer
  \IfBooleanTF{#1}{%
    % *-version => auto-sized braces
    \Bracket*{#3}%
  }{%
    % no star
    \IfNoValueTF{#2}{%
      % no optional argument => normal size braces
      \Bracket{#3}%
    }{%
      % optional argument => \big, \Big, etc.
      \Bracket[#2]{#3}%
    }%
  }%
}}

% Expectation:  \E{X} => E[X],   \E*{X} => auto-sized [ X ]
\NewDocumentCommand{\Expect}{som}{%
\ensuremath{\mathbb{E}
  \IfBooleanTF{#1}{%
    \Bracket*{#3}%
  }{%
    \IfNoValueTF{#2}{%
      \Bracket{#3}%
    }{%
      \Bracket[#2]{#3}%
    }%
  }%
}}

\NewDocumentCommand{\Max}{s o m m}{%
\ensuremath{\max
  \IfBooleanTF{#1}{%
    \Brace*{#3,\, #4}%
  }{%
    \IfNoValueTF{#2}{%
      \Brace{#3,\, #4}%
    }{%
      \Brace[#2]{#3,\, #4}
    }%
  }%
}}

\NewDocumentCommand{\Min}{s o m m}{%
\ensuremath{\min
  \IfBooleanTF{#1}{%
    \Brace*{#3,\, #4}%
  }{%
    \IfNoValueTF{#2}{%
      \Brace{#3,\, #4}%
    }{%
      \Brace[#2]{#3,\, #4}
    }%
  }%
}}

% Prefer to ues \Exp{} for Experiment notation 

% Exponential:  \Exp{x} => exp(x)
% \NewDocumentCommand{\Exp}{som}{%
% \ensuremath{\exp
%   \IfBooleanTF{#1}{%
%     \Paren*{#3}%
%   }{%
%     \IfNoValueTF{#2}{%
%       \Paren{#3}%
%     }{%
%       \Paren[#2]{#3}%
%     }%
%   }%
% }}

% You might want \Log, \Var, \Cov, etc. Similarly:
\NewDocumentCommand{\Log}{som}{%
\ensuremath{\log
  \IfBooleanTF{#1}{%
    \Paren*{#3}%
  }{%
    \IfNoValueTF{#2}{%
      \Paren{#3}%
    }{%
      \Paren[#2]{#3}%
    }%
  }%
}}

\NewDocumentCommand{\Var}{som}{%
\ensuremath{\mathrm{Var}
  \IfBooleanTF{#1}{%
    \Bracket*{#3}%
  }{%
    \IfNoValueTF{#2}{%
      \Bracket{#3}%
    }{%
      \Bracket[#2]{#3}%
    }%
  }%
}}


\NewDocumentCommand{\SetBuilder}{s o m m}{%
\ensuremath{\IfBooleanTF{#1}{%
    % Star version: auto-sized braces
    \Brace*{\,#3 : #4\,}%
  }{%
    % Non-star version
    \IfNoValueTF{#2}{%
      % No optional argument => normal-size braces
      \Brace{\,#3 : #4\,}%
    }{%
      % Optional argument => user-specified size, e.g. \big, \Big, \bigg
      \Brace[#2]{\,#3 : #4\,}%
    }%
  }%
}}

%% Additional Macros

%% Type-setting for the different objects

% Primitives: \prim{x} => denote the cryptographic primitive x
\NewDocumentCommand{\prim}{m}{
\ensuremath{
	\textup{\textsf{#1}}
}
}

% Properties: \prop{x} => denote the property x
\NewDocumentCommand{\prop}{m}{
  \ensuremath{
    \textup{#1}
  }
}

% \alg{x}{y} => Algorithm x of construction y
\NewDocumentCommand{\alg}{smm}{
\ensuremath{
	\IfBooleanTF{#1}{
    \textup{\textsf{#2}}_{#3}
    }{
      \textup{\textsf{#2}}_{\prim{#3}}
	}
}
}

% \Exp{x,A}{z} => z security experiment/game for instance x and adversary A
\NewDocumentCommand{\Exp}{smm}{
\ensuremath{
	\IfBooleanTF{#1}{
      \textup{\textbf{Exp}}_{#2}^{\textup{#3}}
      }{
      \textup{\textbf{Exp}}_{\prim{#2}}^{\textup{#3}}
	}
}
}

% \Adv{x,A}{z} => advantage of adversary A in \Exp{x,A}{z}
\NewDocumentCommand{\Adv}{smmm}{
\ensuremath{
	\IfBooleanTF{#1}{
		\textup{\textbf{Adv}}_{#2}^{\textup{#3}}(#4)
	}{
		\textup{\textbf{Adv}}_{\prim{#2}}^{\textup{#3}}(#4)
	}
}
}

% \impl{x} => denote the implementation x
\NewDocumentCommand{\impl}{m}{
	\textup{\text{#1}}
}

% \todo[x] => x needs to be done here
% \NewDocumentCommand{\todo}{o}{
% 	\ensuremath{
% 		\IfNoValueTF{#1}{
% 			\color{red}
% 			\textbf{ Todo! }
% 			\color{black}
% 		}{
% 			\color{red}
% 			\textbf{ TODO: }\text{#1}
% 			\color{black}
% 		}
% 	}
% }

% \negl[n] => function negligible in n
\NewDocumentCommand{\negl}{so}{
  \ensuremath{
    \IfNoValueTF{#2}{
      \textup{\textsf{negl}}(\lambda)
    }{
      \IfBooleanTF{#1}{
        \textup{\textsf{negl}}\Paren*{#2}
      }{
        \textup{\textsf{negl}}\Paren{#2}
      }
    }
  }
}

% \poly[n] => function polynomial in n
\NewDocumentCommand{\poly}{so}{
  \ensuremath{
    \IfNoValueTF{#2}{
      \textup{\textsf{poly}}(\lambda)
    }{
      \IfBooleanTF{#1}{
        \textup{\textsf{poly}}\Paren*{#2}
      }{
        \textup{\textsf{poly}}\Paren{#2}
      }
    }
  }
}

% Robust combiner
\NewDocumentCommand{\Comb}{smo}{
  \ensuremath{
    \IfNoValueTF{#3}{
      \Pi_{\textup{\textbf{#2}}}
    }{
      \IfBooleanTF{#1}{
        \Pi_{\textup{\textbf{#2}}}\Paren*{#3}
      }{
        \Pi_{\textup{\textbf{#2}}}\Paren{#3}
      }
    }
  }
}

% Robust black-box combiner
\NewDocumentCommand{\BBComb}{mo}{
  \ensuremath{
    \IfNoValueTF{#2}{
      {\Pi_{\prim{#1}}}
    }{
      {\Pi_{\prim{#1}}^{#2}}
    }
  }
}

% Interpretation
\NewDocumentCommand{\inter}{sm}{
  \IfBooleanTF{#1}{
    \Bracket*{\!\Bracket*{#2}\!}
  }{
    \Bracket{\!\Bracket{#2}\!}
  }
}

%% The xspace package is used to automatically handle spacing, regardless of whether the command is followed by punctuation or not.
\newcommand{\onemax}{\textsc{OneMax}\xspace}

%% \eps is quicker to type and the \varepsilon font is more common and nice.
\newcommand{\eps}{\ensuremath\varepsilon}

%% Also set the alternate phi as default, no real reason to use the old one, but one can.
\let\oldphi\phi
\renewcommand{\phi}{\ensuremath{\varphi}}
